% Generated from validated Article Draft Result. Do not hand-edit; revise the structured draft instead.

2026年1月を振り返ると、ここで完成したのは一つの『agentic engineering』という完成形ではない。Codex loopや社内data agentが示したharness/context、外部actionに伴うsecurity boundary、simulationを使うtraining、tool-awareなmodel/API、そしてそれを受け止めるserving runtimeが、別々の場所で同時に立ち上がった月だった。重要なのは、base model以外の層がAIの実用性を決める比重を増したことだ。\autocite{src-0d2abbb0fa62,src-3bc80a3c7870,src-85bc3e4f7b39,src-ae10eb4597fe,src-eac8dfbeb78c,src-fa5410010e0a}


\begin{itemize}
  \item 後続号へつながる観察軸は五つある。第一にharness/context/tool execution、第二にsecurityとhuman control、第三にreal environment依存を下げるtraining environment、第四にtool-native model/API、第五にnew workloadを吸収するserving/runtimeである。これらは1月のEvidenceから確認できる構成要素であり、後続月を読むための座標として使える。
\end{itemize}
\autocite{src-0d2abbb0fa62,src-85bc3e4f7b39,src-ae10eb4597fe,src-d9f9de620203,src-eac8dfbeb78c}

\begin{claimboundary}
この連続線はretrospectiveな読み方であって、2月〜7月の個別event、performance、availabilityを1月の事実へ遡及させるものではない。DynaWebはJanuary original submissionでchronologyを固定し、Qwenのliving lifecycle pageはdated entryとregionを固定し、vLLMは1月20日と29日の二releaseを分けて読む。\autocite{src-3bc80a3c7870,src-85bc3e4f7b39,src-d9f9de620203,src-eac8dfbeb78c}
\end{claimboundary}

\begin{claimboundary}
同様に、OpenAIの二つのagent-system記事は別systemとして帰属を分け、internal data-agentの経験を普遍benchmarkへ一般化しない。safe-URL safeguardはURL-based exfiltrationへのdefense-in-depthでありprompt injection一般の解決ではなく、そのsecurity propertyも独立penetration test済みではない。runtimeのperformance claimもproject-reportedのまま扱う。\autocite{src-0d2abbb0fa62,src-3bc80a3c7870,src-ae10eb4597fe,src-d9f9de620203,src-fa5410010e0a}
\end{claimboundary}

\begin{claimboundary}
training側にも未確定境界が残る。DynaWebのsimulated web world modelはadversarialで動的なlive webとmaterially異なり得る。1月時点で確認できるのは、simulationを使ってtraining costと環境依存を制御する方向が現れたことまでであり、実世界への完全なtransferではない。\autocite{src-85bc3e4f7b39}
\end{claimboundary}

したがって1月の転換点は、Agentが『強いmodelの別名』ではなくなったことにある。実行を支えるharness、context、permission、training environment、tool-enabled API、runtimeがそれぞれ独立した設計課題として姿を現した。次の月々で問うべきなのは、このstackがどの層で統合され、どこに新しいbottleneckとcontrol boundaryが生まれるかである。\autocite{src-0d2abbb0fa62,src-3bc80a3c7870,src-85bc3e4f7b39,src-ae10eb4597fe,src-eac8dfbeb78c}
